\textbf{Phần 1. Trắc nghiệm lựa chọn}
\setcounterpageref{ex}{0}

\begin{ex} %[3D2V14]
Một đàn gà có $36$ con. $\frac{1}{4}$ số con gà trong đàn là bao nhiêu con?
\choice
{$8$ con gà}
{\True $9$ con gà}
{$10$ con gà}
{$12$ con gà}
\loigiai{
$\frac{1}{4}$ số gà là: $36 : 4 = 9$ (con).
\textbf{Đáp án của câu hỏi là B}
}
\end{ex}

\begin{ex} %[3D2V14]
Một rổ có $48$ quả cam. Mẹ đã bán đi $\frac{1}{6}$ số quả cam đó. Hỏi mẹ còn lại bao nhiêu quả cam?
\choice
{\True $40$ quả cam}
{$36$ quả cam}
{$42$ quả cam}
{$38$ quả cam}
\loigiai{
Số cam mẹ đã bán là: $48 : 6 = 8$ (quả).\
Số cam còn lại là: $48 - 8 = 40$ (quả).
\textbf{Đáp án của câu hỏi là A}
}
\end{ex}

\begin{ex} %[3D2V14]
Tìm một số biết rằng $\frac{1}{7}$ của số đó bằng $9$. Số đó là
\choice
{$54$}
{$56$}
{\True $63$}
{$72$}
\loigiai{
Số cần tìm là: $9 \times 7 = 63$.
\textbf{Đáp án của câu hỏi là C}
}
\end{ex}

\begin{ex} %[3D2V14]
Một đoạn thẳng $AB$ dài $32\text{ cm}$. Điểm $M$ thuộc đoạn thẳng $AB$ sao cho đoạn thẳng $AM$ bằng $\frac{1}{4}$ đoạn thẳng $AB$. Độ dài đoạn thẳng $MB$ là
\choice
{$8\text{ cm}$}
{$16\text{ cm}$}
{$20\text{ cm}$}
{\True $24\text{ cm}$}
\loigiai{
Độ dài đoạn thẳng $AM$ là: $32 : 4 = 8\text{ (cm)}$.\
Độ dài đoạn thẳng $MB$ là: $32 - 8 = 24\text{ (cm)}$.
\textbf{Đáp án của câu hỏi là D}
}
\end{ex}

\begin{ex} %[3D2V14]
Lớp 3A có $40$ học sinh. Trong đó có $\frac{1}{5}$ số học sinh là học sinh xuất sắc, $\frac{1}{4}$ số học sinh là học sinh tiêu biểu. Hỏi lớp 3A có tất cả bao nhiêu bạn đạt danh hiệu khen thưởng trên?
\choice
{$16$ bạn}
{\True $18$ bạn}
{$20$ bạn}
{$15$ bạn}
\loigiai{
Số học sinh xuất sắc là: $40 : 5 = 8$ (bạn).\
Số học sinh tiêu biểu là: $40 : 4 = 10$ (bạn).\
Tổng số bạn được khen thưởng là: $8 + 10 = 18$ (bạn).
\textbf{Đáp án của câu hỏi là B}
}
\end{ex}

\begin{ex} %[3D2V14]
So sánh $\frac{1}{6}$ của $54\text{ kg}$ và $\frac{1}{8}$ của $64\text{ kg}$:
\choice
{\True $\frac{1}{6}$ của $54\text{ kg}$ lớn hơn}
{$\frac{1}{8}$ của $64\text{ kg}$ lớn hơn}
{Hai giá trị bằng nhau}
{Không so sánh được}
\loigiai{
$\frac{1}{6}$ của $54\text{ kg}$ là: $54 : 6 = 9\text{ kg}$.\
$\frac{1}{8}$ của $64\text{ kg}$ là: $64 : 8 = 8\text{ kg}$.\
Vì $9\text{ kg} > 8\text{ kg}$ nên $\frac{1}{6}$ của $54\text{ kg}$ lớn hơn.
\textbf{Đáp án của câu hỏi là A}
}
\end{ex}

\begin{ex} %[3D2V14]
Một ngày có $24$ giờ. Hỏi $\frac{1}{3}$ ngày có bao nhiêu giờ?
\choice
{$6$ giờ}
{$12$ giờ}
{\True $8$ giờ}
{$4$ giờ}
\loigiai{
Số giờ trong $\frac{1}{3}$ ngày là: $24 : 3 = 8$ giờ.
\textbf{Đáp án của câu hỏi là C}
}
\end{ex}

\begin{ex} %[3D2V14]
Một cửa hàng có $56\text{ m}$ vải hoa. Cửa hàng đã bán được $\frac{1}{7}$ số vải hoa đó và bán thêm $12\text{ m}$ nữa. Hỏi cửa hàng đã bán tất cả bao nhiêu mét vải hoa?
\choice
{$18\text{ m}$}
{$16\text{ m}$}
{$22\text{ m}$}
{\True $20\text{ m}$}
\loigiai{
$\frac{1}{7}$ số vải hoa là: $56 : 7 = 8\text{ (m)}$.\
Tổng số vải đã bán là: $8 + 12 = 20\text{ (m)}$.
\textbf{Đáp án của câu hỏi là D}
}
\end{ex}

\begin{ex} %[3D2V14]
Bao gạo thứ nhất nặng $45\text{ kg}$. Bao gạo thứ hai nặng bằng $\frac{1}{5}$ bao gạo thứ nhất. Hỏi cả hai bao nặng bao nhiêu ki-lô-gam?
\choice
{$50\text{ kg}$}
{\True $54\text{ kg}$}
{$48\text{ kg}$}
{$56\text{ kg}$}
\loigiai{
Bao thứ hai nặng: $45 : 5 = 9\text{ kg}$.\
Cả hai bao nặng: $45 + 9 = 54\text{ kg}$.
\textbf{Đáp án của câu hỏi là B}
}
\end{ex}

\begin{ex} %[3D2V14]
Hình vuông có chu vi bằng $36\text{ cm}$. Độ dài một cạnh của hình vuông là
\choice
{\True $9\text{ cm}$}
{$6\text{ cm}$}
{$8\text{ cm}$}
{$12\text{ cm}$}
\loigiai{
Cạnh hình vuông bằng $\frac{1}{4}$ chu vi: $36 : 4 = 9\text{ cm}$.
\textbf{Đáp án của câu hỏi là A}
}
\end{ex}

\begin{ex} %[3D2V14]
Một tổ công nhân nhận sửa một đoạn đường dài $72\text{ m}$. Ngày thứ nhất sửa được $\frac{1}{9}$ đoạn đường, ngày thứ hai sửa được $\frac{1}{8}$ đoạn đường. Hỏi sau hai ngày tổ công nhân đã sửa được bao nhiêu mét đường?
\choice
{$15\text{ m}$}
{$16\text{ m}$}
{\True $17\text{ m}$}
{$18\text{ m}$}
\loigiai{
Ngày 1 sửa: $72 : 9 = 8\text{ m}$.\
Ngày 2 sửa: $72 : 8 = 9\text{ m}$.\
Cả hai ngày sửa: $8 + 9 = 17\text{ m}$.
\textbf{Đáp án của câu hỏi là C}
}
\end{ex}

\begin{ex} %[3D2V14]
Một thùng dầu chứa $54\text{ lít}$. Người ta lấy ra $\frac{1}{6}$ số dầu rồi lại lấy thêm $5\text{ lít}$ nữa. Số lít dầu còn lại trong thùng là
\choice
{$35\text{ lít}$}
{$45\text{ lít}$}
{$42\text{ lít}$}
{\True $40\text{ lít}$}
\loigiai{
Lần thứ nhất lấy: $54 : 6 = 9\text{ lít}$.\
Cả hai lần lấy: $9 + 5 = 14\text{ lít}$.\
Còn lại: $54 - 14 = 40\text{ lít}$.
\textbf{Đáp án của câu hỏi là D}
}
\end{ex}

\textbf{Phần 2. Trắc nghiệm đúng sai}
\setcounterpageref{ex}{0}

\begin{ex} %[3D2V14]
Xét tính đúng sai của các nhận định về phân số một phần mấy:
\choiceTF
{\True Để tìm một phần ba của số $A$, ta lấy $A$ chia cho 3.}
{Một phần năm của 45 lớn hơn một phần tư của 40.}
{\True Nếu $\frac{1}{8}$ số bi của Nam là 6 viên thì Nam có tất cả 48 viên bi.}
{\True Một giờ có 60 phút thì $\frac{1}{6}$ giờ tương ứng với 10 phút.}
\loigiai{
\begin{itemize}
    \item Ý a) Đúng: Quy tắc tìm một phần mấy của một số.
    \item Ý b) Sai: $\frac{1}{5}$ của $45$ là $9$, $\frac{1}{4}$ của $40$ là $10$. Vì $9 < 10$ nên nhận định Sai.
    \item Ý c) Đúng: Số bi của Nam là: $6 \times 8 = 48$ viên.
    \item Ý d) Đúng: $60 : 6 = 10$ phút.
\end{itemize}
\begin{itemize}
    \item a) \textbf{Đúng}.
    \item b) \textbf{Sai}.
    \item c) \textbf{Đúng}.
    \item d) \textbf{Đúng}.
\end{itemize}
}
\end{ex}

\begin{ex} %[3D2V14]
Một trang trại có 72 con cừu. Bác nông dân đã bán đi $\frac{1}{8}$ đàn cừu:
\choiceTF
{\True Bác nông dân đã bán đi 9 con cừu.}
{\True Đàn cừu còn lại 63 con.}
{Số cừu còn lại gấp 6 lần số cừu đã bán.}
{\True Nếu bác bán thêm 7 con cừu nữa thì số cừu còn lại là 56 con.}
\loigiai{
\begin{itemize}
    \item Ý a) Đúng: $72 : 8 = 9$ con.
    \item Ý b) Đúng: $72 - 9 = 63$ con.
    \item Ý c) Sai: $63 : 9 = 7$ lần (không phải 6 lần).
    \item Ý d) Đúng: $63 - 7 = 56$ con.
\end{itemize}
\begin{itemize}
    \item a) \textbf{Đúng}.
    \item b) \textbf{Đúng}.
    \item c) \textbf{Sai}.
    \item d) \textbf{Đúng}.
\end{itemize}
}
\end{ex}

\textbf{Phần 3. Trả lời ngắn}
\setcounterpageref{ex}{0}

\begin{ex} %[3D2V14]
Tìm $\frac{1}{7}$ của $63\text{ kg}$.
\shortans{9}
\loigiai{
\textbf{Đáp số:} $9$\\
$63 : 7 = 9\text{ kg}$.
}
\end{ex}

\begin{ex} %[3D2V14]
Tìm một số biết rằng $\frac{1}{6}$ của số đó bằng 8.
\shortans{48}
\loigiai{
\textbf{Đáp số:} $48$\\
$8 \times 6 = 48$.
}
\end{ex}

\begin{ex} %[3D2V14]
Một mảnh vườn hình chữ nhật có chu vi bằng $40\text{ m}$. Nửa chu vi mảnh vườn đó là bao nhiêu mét?
\shortans{20}
\loigiai{
\textbf{Đáp số:} $20$\\
Nửa chu vi bằng $\frac{1}{2}$ chu vi: $40 : 2 = 20\text{ m}$.
}
\end{ex}

\begin{ex} %[3D2V14]
Lớp 3B có 36 học sinh. Số học sinh nam bằng $\frac{1}{4}$ số học sinh cả lớp. Hỏi lớp 3B có bao nhiêu học sinh nữ?
\shortans{27}
\loigiai{
\textbf{Đáp số:} $27$\\
Số nam: $36 : 4 = 9$ bạn. Số nữ: $36 - 9 = 27$ bạn.
}
\end{ex}

\begin{ex} %[3D2V14]
Một bao gạo nặng 54 kg. Người ta đã dùng hết $\frac{1}{9}$ bao gạo đó. Hỏi trong bao còn lại bao nhiêu ki-lô-gam gạo?
\shortans{48}
\loigiai{
\textbf{Đáp số:} $48$\\
Đã dùng: $54 : 9 = 6\text{ kg}$. Còn lại: $54 - 6 = 48\text{ kg}$.
}
\end{ex}

\begin{ex} %[3D2V14]
Một cuộn dây dài 45 m. Người ta cắt lấy $\frac{1}{5}$ cuộn dây. Hỏi đoạn dây lấy ra dài bao nhiêu mét?
\shortans{9}
\loigiai{
\textbf{Đáp số:} $9$\\
$45 : 5 = 9\text{ m}$.
}
\end{ex}
